\chapter[Sermon 65. The Faithful Assuredly and Straightway Flit from the Corporal Death unto Life Everlasting.]{The Faithful Assuredly and Straightway Flit from the Corporal Death unto Life Everlasting.}
\label{sermon:065}
\markright{Sermon 65. The Faithful Assuredly and Straightway Flit from the ...}

And I heard a voice from heaven, saying to me: Write, “Blessed are the dead who hereafter die in the Lord.” Indeed, the spirit says that they rest from their labors. But their works follow them.

Although he has spoken on many occasions about the state of souls in the world beyond, and about the blessedness of the faithful who are killed for the sake of religion, it was chiefly necessary to address this matter here. For I asked, how many must be killed by the beast. Now, lest they, fearing death, should choose rather to worship the beast than to be slain, and lest, perhaps, having lost this life, there be no other life to be hoped for in the world to come, he diligently and certainly treats of the state of souls, and of the felicity and blessedness of souls, which as soon as they die, they attain, assuredly and immediately departing from this world into everlasting life. But those who know these things, and have conceived them through true faith, and understand that they will undoubtedly pass from bodily death into blessed life, cannot but more boldly despise this present life.

And this wholesome doctrine is comprised in three points. First, he shows its certainty; secondly, he declares what it is; last, he sets forth and illuminates it by circumstances. At the first, verily, he seems to allude to the manner, customarily received by all nations, that such things as they would have thought to be certain and undoubted, they would also commit to writing to leave them unto posterity. But the certainty and verity or authority of the thing is esteemed by authors, who first have dispatched any matters among themselves, and after have caused the same to be put in writing. At this present, therefore, is God shown to be author. For John says: and I heard a voice from heaven. And by and by adds: the spirit says. Therefore, there is no doubt but that the Son of God himself has spoken and revealed these things.

\index{Antichrist}\index{John, Saint}\index{Papacy}\index{Rome}For him he saw at the beginning of this revelation; after he sees diverse kinds of angels, but he sees not Christ speaking to him. But he hears now his voice from heaven; he hears the spirit speaking, by whom the Lord said, whilst he was yet conversant in earth with his disciples, that he would treat and speak all things in the church. Let us believe therefore that the words which are here recited, by Christ’s doing, are a celestial oracle, certain and true, whereof we ought not to doubt. And John the Apostle and Evangelist is commanded to write the sayings of Christ from the heavenly seat. Which thing he doth; and so at Christ’s commandment sends them unto all posterity, unto us also and to our offspring even to the world’s end. But if writings written by the chancellors or secretaries of kings and princes, being notable men, deserve credit, we may much more justly and rightly believe this writing, which the Son of God indites from heaven; and that beloved disciple of Christ, the apostle and Evangelist John writes. You had once a confidence in the Popes’ bulls (they may well be called bulls, since they are more vain than bulls or blabbers in the water), sent from the See of Rome, wherein you, as one assured, did put full trust to have remission of sins and blessed life. And shall you not now be accounted mad and out of your wit, in case you will not believe this heavenly writing? That other was indicted by the spirit of Antichrist, by the Pope, the man of sin, and child of perdition; and written of some deceiver infected with Simony and sacrilege, which in life and manners was filthiness itself. But in John is nothing but cleanness, purity, and integrity; and the very Son of God which prescribes these things to John, is the very verity and life, the light of the world and lord of heaven and earth, of life and death. See then how safely you may lay yourself to this heavenly writing, which here is offered and given to you freely. You need not to disburse for the same one farthing. The Pope instituted in the church buying and selling and devilish bargaining about indulgences and other things, which were plain deceipts and illusions, plain mockeries, and open blasphemies, and therefore accursed forever; as Peter also pronounces in the eighth chapter of the Acts. God himself dissuades all men from such tromperies, and bargains wicked and vain, in the fifty-fifth chapter of Isaiah, where he promises again that he will give to the godly all plenty of all good things.

And now let us hear what the writing is, and what Saint John is commanded from heaven to put in writing. It is a short sentence, as also in many places, the wisdom of God comprehends in few words the true sum of blessedness: so providing for our infirmity, that we need not to complain that the doctrine were too long, which we with our slender understanding are not able to attain to. The Lord therefore pronounces them to be blessed who die in the Lord; then we must see what he understands by blessedness, and who they are that die in the Lord. Blessedness is that high felicity which happens to the faithful in another world, in which we shall see God himself as he is, and have the fruition of him unto a joyful and never-wearisome fullness. We shall live in the same with all the Saints forever, and shall have joys that cannot be expressed with the tongues of men. Of which more shall follow afterward. They shall rest from their labors, and more plentifully in the twenty-first chapter. And they die in the Lord who, by faith grafted in Christ, are joined to him alone, depend wholly upon him, only regard him, and desire nothing else but him alone. For they are said to live in Christ, in whom Christ lives by faith; they that live in Christ do frame their whole life after the will of Christ. And they die in the Lord chiefly and before all, who for the confession of the Lord’s faith, suffer death, and offer themselves to torments. And not they alone, but those also, who although they die by the sword of the persecutors, yet die when the Lord calls them in the true Christian faith. For these are also blessed, as the Lord in Saint John says, “Verily, verily, I say unto you, if any man keep my word, he shall not see death forever.” However, they do not die in the Lord who either deny God that they might not be slain, or trust to their own merits and intercessions of Saints, or to other men’s works, be they monks, friars, or mass-serving priests, and so depart out of this life, thinking that they shall be helped by other men’s works. To be brief, the truth of the Lord pronounces them all blessed and fortunate who depart out of this world in true faith.

\index{Erasmus of Rotterdam}\index{Arethas of Caesarea}\index{Resurrection}\index{Judgment, Last}Finally, the Lord himself adds a notable declaration of this brief statement. For he sets forth the circumstances of the time and the manner of the blessedness. For it is wont to be demanded, what time salvation and felicity happens to the dead? whether immediately, or after a time? That is, do our souls depart immediately after the death of the body to the blessed places? Or are they intercepted for a certain time, so that they might be purged in purgatory before they enter into heaven? Or are they held in sleep, awaiting the resurrection of the bodies, so that they might then awaken and together with their bodies enter into heaven? To all of which things the heavenly oracle answers immediately, saying that the same felicity comes to souls immediately. In the Latin copies, this place is pointed thus: “Blessed are the dead who die in the Lord.” Immediately now says the Spirit, that they may rest from their labors. In like manner reads the Spanish or Complutensian copy. But Arethas and the Greek copies, and also the exemplar of Paris are pointed so that [illegible] should be the end of the sentence, as Erasmus notes. After follows [illegible], which is, “verily, certainly,” says the Spirit. The sense is therefore, that the faithful, being dead, shall straightway and immediately achieve salvation. For the word which Saint John uses signifies from the very instant, from that hour immediately, incontinently. This allows no space, but expresses that which we are wont to note by the Dutch phrase. Being admonished therefore by a divine oracle and confirmed by a writ brought from heaven, let us all be assured that the souls of all the faithful depart from bodily death into everlasting life. These things are confirmed and made plain also by other passages of Scripture innumerable. I will choose out only a certain few, and those also the testimonies of our Savior, who is the light of the world and the word of life. In chapter 3 of Saint John, he says expressly that the faithful are delivered from death by his cross, as in times past, by the sight of the brazen serpent, the Israelites were delivered from the deadly sting of venomous poison. And it is plain that they were delivered immediately and most fully. In John 5, the same says, “I have passed from death to life.” Let this passage be weighed diligently, and it shall appear that it alone suffices in this matter. In John 6, he says openly, “I will raise him on the last day.” But he raises not the bodies only at the last judgment, but in every man’s last day, that is, in the death of every one, he preserves the souls, that they should not perish or be tormented, and so forth. We have in the gospel examples most clear: namely, of Lazarus the beggar, who was immediately after his death carried up by angels into the bosom of Abraham; and of the thief, who heard the Lord say, “Today you will be with me in Paradise”; and of Stephen saying, “Lord Jesus, receive my spirit”; but especially of our Savior, saying on the cross, “Father, into your hands I commend my spirit,” and so forth.

By these things, whatever the monastic and Antichristian doctrine has built—regarding purgatory, indulgences, and the miserable state of souls in the world to come—is completely overturned. From this, they made a most shameful profit. Those who believe that souls are mortal are also refuted, as well as those who believe that souls sleep in another world, where they cannot even sleep here in this state of weakness. Therefore, you will say it is madness to think that souls sleep, being freed from the burden of the body.

Concerning the manner of the blessedness of the saints, they rest from their labors. Therefore, salvation is a most joyful tranquility. All diseases, sicknesses, griefs, anxieties, sorrow, famine, thirst, cold—briefly, all things that vex or trouble people—are removed. Rest and tranquility, joy and blessing take their place. And since the dead rest from their labors, who can believe that they are afflicted with torments? But lest any man should doubt even a little, he adds a confirmation, [illegible], yes, or certainly, verily, the spirit says, the dead shall be quiet from all their griefs. Let no man therefore doubt.

And he adds another thing, that the works of the saints follow, that is to say, after the saints have departed from this world, they are rewarded in another world, if they have done anything well, if they have suffered hard things. For a reward is prepared for virtues. The saints hope for and receive this reward without boasting of their own merit, and not in contempt of the merit of Christ. For they acknowledge that God in his saints crowns his own gifts. And this is spoken of the reward of works for the consolation of those who suffer many things in this world. So the Lord said in the Gospel: “Your reward is plentiful in heaven.” And the Apostle affirms everywhere that rewards are prepared for those who are crucified here with Christ. And here let us mark diligently that these things are spoken also of the spirit of Christ under the guise of another. For the world despises religious persons and those who suffer for religion, and objects that they lose their labor and cost. Contrarywise, the spirit by another [?vouchsafes] avouches that a reward is prepared for virtue.

Let us also note this: their works, and not those of other men, are what we should follow, and they are not sent after by others. Let no one deceive himself, let no one think that after his death, priests will send him into purgatory a portion of other men’s merits. Those are not good works which are done by priests and friars apart from and against God’s word, but provocations of God’s wrath. Those are not excluded from the kingdom of God in the Gospel who go to others to buy oil. Scripture says elsewhere: Let us do good while we have time, for a time will come when no one can work. Let us therefore watch, and in deed do good works through faith.
